\chapter{Proof Program and Mechanized Foundations}
\label{ch:proof-program}

\textbf{Status: proof contracts and proof strategies are integrated here as an
expository projection. All physics proof packages remain \texttt{proposed}.
Only \texttt{PRF-05.01} has been checked by one pinned Lean backend; no contract
is cross-checked, and no later reduction theorem is represented as complete or
reviewed.}

The project maintains a repository-owned proof workspace and a separate
prover-neutral theorem catalog. Those records own proof decomposition and
status. This chapter connects them to the scientific narrative without
promoting a draft theorem, proof sketch, numerical check, or manuscript label
into stronger evidence.

\section{Why a separate proof program is needed}

The computational hierarchy contains claims of several different kinds. Some
are elementary finite-dimensional identities, such as covariance under a
unitary coordinate change. Some require analytical hypotheses, such as a
spectral gap, compact parameter domain, resolvent bound, scale separation, or
asymptotic locality. Others are numerical or scientific statements that cannot
be established by proof alone.

A successful calculation may suggest a theorem but does not prove it. A formal
proof may establish an implication from assumptions while leaving the physical
adequacy of those assumptions unvalidated. The proof program therefore records

\begin{itemize}
  \item state spaces, domains, codomains, and group actions;
  \item assumptions needed by each result;
  \item analytical dependencies among results;
  \item mechanized theorem contracts and backend status;
  \item nonclaims and failure boundaries; and
  \item numerical or scientific evidence that remains separate from proof.
\end{itemize}

The proof workspace is subordinate to the versioned physical and numerical
specifications. It cannot repair an unresolved scientific convention by
choosing a mathematically convenient one.

\section{Proof-package graph}

The physical proof program is divided into six packages:

\begin{description}
  \item[\texttt{PRF-00}: Foundations.] State spaces, Bloch-fiber
    correspondence, and representation and reduction maps.
  \item[\texttt{PRF-05}: Mechanized operator lemmas.] Nine frozen
    finite-dimensional contracts targeted independently in Lean, Isabelle, and
    Rocq.
  \item[\texttt{PRF-10}: Operator alignment.] TB-anchored identification,
    gauge actions, and aligned pristine--doped subtraction.
  \item[\texttt{PRF-20}: Reduction.] Gauge-constrained locality,
    excluded-space reduction, atomistic-to-continuum consistency, and path
    commutativity.
  \item[\texttt{PRF-30}: Bounds.] Crossover, operator-to-observable, and
    continuum-fitting statements.
  \item[\texttt{PRF-40}: Compatibility.] Spectral/operator admissibility,
    model-class distance, and certified incompatibility.
\end{description}

Their dependency structure is a partial order rather than a single pipeline:
foundations support every later package; mechanized elementary lemmas support
applicable alignment, reduction, and compatibility claims; alignment and
reduction support the proposed bounds; and compatibility also depends on
explicit model classes and admissible gauges.

The separate \texttt{PRF-90} package concerns agentic-workflow properties such
as authorization safety and replay determinism. It is intentionally outside the
physics proof chain. A workflow theorem would not establish a physical or
numerical theorem.

\section{Status vocabulary}

The proof registry distinguishes \texttt{proposed}, \texttt{in development},
\texttt{blocked}, \texttt{complete, unreviewed}, \texttt{reviewed},
\texttt{numerically checked}, and \texttt{retired}. Mathematical review,
numerical verification, scientific validation, and human publication acceptance
remain distinct.

A frozen mechanization contract fixes the current theorem target; it does not
mean that the theorem has been encoded or proved. A backend status of
\texttt{checked} means that the identified checker accepted the identified
encoding under the pinned toolchain. Cross-checking additionally requires all
three independent backends and semantic conformance review.

The derivations below explain selected contracts. They do not modify the owning
registry status.

\section{Finite-dimensional retained frames}

Let $V\in\mathbb C^{D\times M}$ have orthonormal columns,
\begin{equation}
  V^{\dagger}V=I_M,
\end{equation}
and let $G\in\mathbb C^{M\times M}$ be unitary. The columns of $V$ represent an
ordered orthonormal frame for an $M$-dimensional retained subspace of an ambient
$D$-dimensional space.

\begin{proposition}[\texttt{PRF-05.01}: projector invariance]
\label{prop:projector-invariance}
For $P=VV^{\dagger}$ and $V'=VG$, the rotated frame induces the same ambient
orthogonal projector:
\begin{equation}
  V'V'^{\dagger}=P.
\end{equation}
\end{proposition}

\begin{proof}[Expository derivation]
Using $(VG)^{\dagger}=G^{\dagger}V^{\dagger}$ and
$GG^{\dagger}=I_M$,
\begin{equation}
  (VG)(VG)^{\dagger}
  =VGG^{\dagger}V^{\dagger}
  =VV^{\dagger}.
\end{equation}
The condition $V^{\dagger}V=I_M$ is needed to interpret $VV^{\dagger}$ as the
orthogonal projector, although the displayed equality uses only the unitarity
of $G$.
\end{proof}

This is the sole frozen contract currently encoded and checked in Lean. The
pinned trial uses Lean~4 and mathlib \texttt{v4.33.0}; the project source contains
no \texttt{sorry} or \texttt{admit}. The check establishes acceptance of that
encoding, not physical admissibility of the frame, smoothness in $\mathbf k$,
or conformance with Isabelle or Rocq.

\begin{proposition}[\texttt{PRF-05.02}: covariance of compression]
\label{prop:compression-covariance}
For an ambient operator $A\in\mathbb C^{D\times D}$, define
\begin{equation}
  H=V^{\dagger}AV,
  \qquad
  H'=(VG)^{\dagger}A(VG).
\end{equation}
Then
\begin{equation}
  H'=G^{\dagger}HG.
\end{equation}
\end{proposition}

\begin{proof}[Expository derivation]
Associativity and the adjoint product rule give
\begin{equation}
  (VG)^{\dagger}A(VG)
  =G^{\dagger}(V^{\dagger}AV)G.
\end{equation}
No Hermiticity assumption on $A$ is needed for this coordinate identity.
\end{proof}

The contract does not assert that $A$ preserves the retained subspace or that
the compressed spectrum equals the full ambient spectrum.

\section{Identification pullback and aligned subtraction}

Let $H_b\in\mathbb C^{M_b\times M_b}$ and
$H_d\in\mathbb C^{M_d\times M_d}$ represent pristine and doped retained
operators. Let
\begin{equation}
  U:\mathbb C^{M_b}\longrightarrow\mathbb C^{M_d}
\end{equation}
be an identification map. Under independent unitary frame changes $G_b$ and
$G_d$, define
\begin{equation}
  H_b'=G_b^{\dagger}H_bG_b,
  \qquad
  H_d'=G_d^{\dagger}H_dG_d,
  \qquad
  U'=G_d^{\dagger}UG_b.
\end{equation}

\begin{proposition}[\texttt{PRF-05.03}: pullback covariance]
\label{prop:pullback-covariance}
The doped operator pulled into pristine coordinates transforms as
\begin{equation}
  U'^{\dagger}H_d'U'
  =
  G_b^{\dagger}(U^{\dagger}H_dU)G_b.
\end{equation}
\end{proposition}

\begin{proof}[Expository derivation]
Insert the primed definitions and cancel
$G_dG_d^{\dagger}=I_{M_d}$ on both sides of $H_d$:
\begin{align}
  U'^{\dagger}H_d'U'
  &=G_b^{\dagger}U^{\dagger}G_d
    G_d^{\dagger}H_dG_d
    G_d^{\dagger}UG_b\\
  &=G_b^{\dagger}U^{\dagger}H_dUG_b.
\end{align}
Unitarity of $U$ is not required for this covariance identity.
\end{proof}

Two common-space differences must be distinguished:
\begin{align}
  \Delta H_{d\to b}
  &=U^{\dagger}H_dU-H_b,
  &\Delta H_{d\to b}&:\mathbb C^{M_b}\to\mathbb C^{M_b},\\
  \Delta H_{b\to d}
  &=H_d-UH_bU^{\dagger},
  &\Delta H_{b\to d}&:\mathbb C^{M_d}\to\mathbb C^{M_d}.
\end{align}

\begin{proposition}[\texttt{PRF-05.04}: aligned-difference covariance and equivalence]
\label{prop:aligned-difference-equivalence}
Each difference transforms covariantly in its own common space. If $U$ is a
unitary identification, then
\begin{equation}
  \Delta H_{b\to d}
  =U\Delta H_{d\to b}U^{\dagger},
  \qquad
  \Delta H_{d\to b}
  =U^{\dagger}\Delta H_{b\to d}U.
\end{equation}
\end{proposition}

\begin{proof}[Expository derivation]
Covariance follows by applying Proposition~\ref{prop:pullback-covariance} to the
pullback difference and its pushforward analogue to the doped-space difference.
For unitary $U$,
\begin{align}
  U\Delta H_{d\to b}U^{\dagger}
  &=U(U^{\dagger}H_dU-H_b)U^{\dagger}\\
  &=H_d-UH_bU^{\dagger}
  =\Delta H_{b\to d}.
\end{align}
The inverse relation follows similarly.
\end{proof}

This algebra does not establish that a physically adequate $U$ exists, that it
is smooth or symmetry compatible, or that the scalar energy zeros of $H_b$ and
$H_d$ agree. Those remain prerequisites before either difference is called an
impurity operator.

\section{Invariant norms and gauge-dependent locality}

For a unitary $G$, the Frobenius and induced Euclidean operator norms satisfy
\begin{equation}
  \|G^{\dagger}AG\|_{\mathrm F}=\|A\|_{\mathrm F},
  \qquad
  \|G^{\dagger}AG\|_2=\|A\|_2.
\end{equation}
These are the frozen \texttt{PRF-05.05a} and \texttt{PRF-05.05b} targets. They
support gauge-invariant residuals only when the compared paths transform
conformantly, as required by \texttt{PRF-05.06}.

Spatial blocks and shell decompositions require a narrower statement. A general
unitary rotation can mix orbitals associated with different sites or cells and
therefore change which entries are classified as local or exterior. The frozen
\texttt{PRF-05.07} target is a two-point counterexample showing that gauge
transformation and real-space truncation need not commute. The
\texttt{PRF-05.08} target instead restricts the gauge to declared lattice-local
constant rotations under which shell Frobenius diagnostics remain invariant.

Thus a global unitarily invariant norm and a spatially resolved residual have
different admissible gauge groups. A proof about the first cannot be reused for
the second without the additional locality assumptions.

\section{TB-anchored identification}

Suppose $\widetilde V_b(\mathbf k)$ and
$\widetilde V_d(\mathbf k)$ are orthonormal frames of equal rank after both have
been related to a declared tight-binding anchor. The candidate fiberwise map is
\begin{equation}
  U_d(\mathbf k)
  =
  \widetilde V_d(\mathbf k)
  \widetilde V_b(\mathbf k)^{\dagger}.
\end{equation}
On the corresponding retained spaces, this formula gives an isometric unitary
identification under the stated orthonormality and equal-rank assumptions.

The substantial \texttt{PRF-10.01} obligations are not the elementary product
identity. They concern existence of suitable anchored frames, rank and
conditioning, global smoothness and periodicity, crystal symmetry, and physical
adequacy. \texttt{PRF-10.02} then owns equivariance of the construction, while
\texttt{PRF-10.03} owns the common-space subtraction and its covariance. All
three remain proposed.

\section{Excluded-space reduction}

Let $P+Q=I$ be complementary orthogonal projectors and write the block
components of a self-adjoint operator $H$ as $H_{PP}=PHP$,
$H_{PQ}=PHQ$, and so on. When a real energy $E$ belongs to the resolvent set of
$H_{QQ}$, eliminating the $Q$ component gives the energy-dependent Feshbach
effective operator~\cite{feshbach1958,feshbach1962}
\begin{equation}
  H_{\mathrm{eff}}(E)
  =
  H_{PP}+H_{PQ}(E-H_{QQ})^{-1}H_{QP}.
\end{equation}
If
\begin{equation}
  \Delta_Q=\operatorname{dist}(E,\sigma(H_{QQ}))>0,
\end{equation}
then the self-adjoint resolvent estimate suggests the proposed bound
\begin{equation}
  \left\|
    H_{PQ}(E-H_{QQ})^{-1}H_{QP}
  \right\|
  \leq
  \frac{\|H_{PQ}\|\,\|H_{QP}\|}{\Delta_Q}.
\end{equation}
For $H_{QP}=H_{PQ}^{\dagger}$ this becomes
$\|H_{PQ}\|^2/\Delta_Q$.

The \texttt{PRF-20.02} obligation must keep this real off-spectrum setting
separate from open-system or analytically continued resonance formulations,
where non-Hermiticity and widths require additional theory. No small correction
is inferred without both a coupling estimate and resolvent distance.

\section{Crossover and operator-to-observable bounds}

Let $D_d=\Delta H_d-V_{\mathrm{cont},d}$ be an aligned atomistic-minus-continuum
discrepancy in a common space, and let $P_{>R}$ project onto a declared exterior
region. Define
\begin{equation}
  \eta_d(R)
  =
  \|P_{>R}D_dP_{>R}\|.
\end{equation}
When the exterior subspaces are nested, compression by a smaller exterior
projector cannot increase the operator norm. A crossover criterion may then be
defined by
\begin{equation}
  r_{c,d}(\tau)
  =
  \inf\{R:\eta_d(R)\leq\tau\}.
\end{equation}

The \texttt{PRF-30.01} obligation must distinguish nonemptiness of the set,
existence of the infimum, attainment by a tested radius, numerical estimation,
and physical interpretation. The assumption
$\eta_d(R)\to0$ as $R\to\infty$ is an asymptotic-locality hypothesis, not a
consequence established by general Kohn--Sham DFT.

For self-adjoint atomistic and reduced operators with perturbation
$E=H_{\mathrm{atom}}-H_{\mathrm{red}}$, standard perturbation theory motivates
spectral-distance bounds controlled by $\|E\|$. Invariant-subspace rotation can
be controlled by a Davis--Kahan-type estimate when the target cluster is
separated by a declared spectral gap~\cite{davis1970}. The proposed
\texttt{PRF-30.02} result must identify the finite or infinite setting, norm,
spectral window, gap, state identification, and degeneracy treatment. A
worst-case bound is not assumed sharp, and a small spectral error does not imply
a small operator error.

\section{Continuum fitting and identifiability}

For continuum parameters $\boldsymbol\theta\in\Theta$, define a declared
exterior objective
\begin{equation}
  J_R(\boldsymbol\theta)
  =
  \left\|
    P_{>R}
    [\Delta H_d-V_{\mathrm{cont}}(\boldsymbol\theta)]
    P_{>R}
  \right\|.
\end{equation}
If $\Theta$ is nonempty and compact and $J_R$ is continuous, the extreme-value
theorem gives existence of a minimizer. The proposed \texttt{PRF-30.03}
program separates this existence result from uniqueness, conditioning,
stability, and physical identifiability. An optimizer output supplies none of
those stronger properties automatically.

\section{Compatibility and certified incompatibility}

Let $\mathfrak A_E$ and $\mathfrak A_H$ be nonempty spectral- and
operator-admissible subsets of a declared metric model space. If both sets are
compact, their distance
\begin{equation}
  \delta^{\ast}
  =
  \min_{H_E\in\mathfrak A_E,\,H_H\in\mathfrak A_H}
  d(H_E,H_H)
\end{equation}
is attained. If the compact sets are disjoint, the distance is positive. This
standard existence argument is only part of \texttt{PRF-40.01}: the scientific
application must also establish closure or compactness, continuity of the
losses, and a reproducible admissible gauge group.

A failed local optimizer does not prove that the sets are disjoint or that
$\delta^{\ast}>0$. The proposed \texttt{PRF-40.03} certification obligation
requires a valid global lower bound, for example from analytic estimates,
interval methods, branch-and-bound, exhaustive certified reduction, or an
applicable validated convex relaxation. \texttt{PRF-40.04} then asks what can be
inferred when a nested model class reduces the residual; a smaller residual does
not uniquely identify the physical cause of the improvement.

\section{Reduction-path commutativity}

For aligned operands in one common space, let $\mathcal R$ be a reduction map.
The path defect is
\begin{equation}
  \varepsilon_{\mathrm{path}}
  =
  \left\|
    \mathcal R(H_d-H_b^{\mathrm{aligned}})
    -
    [\mathcal R(H_d)-\mathcal R(H_b^{\mathrm{aligned}})]
  \right\|.
\end{equation}
If the same linear map acts on both operands in the same coordinates, exact
commutativity follows from linearity. This elementary fact does not apply
automatically to nonlinear fitting, separately selected model classes,
gauge-dependent truncation, or independently chosen alignment maps.

The \texttt{PRF-20.04} program distinguishes gauge equivariance from path
commutativity and seeks bounds when exact equality fails. The frozen
\texttt{PRF-05.06} contract addresses gauge invariance of a residual only after
the two paths are already conformant.

\section{Mechanization status}

The current prover-neutral \texttt{PRF-05} target set contains nine frozen
contracts:

\begin{center}
\small
\begin{tabular}{@{}lll@{}}
\toprule
Contract & Subject & Current backend evidence \\
\midrule
\texttt{05.01} & Projector invariance & Lean checked \\
\texttt{05.02} & Compression covariance & Unencoded \\
\texttt{05.03} & Pullback covariance & Unencoded \\
\texttt{05.04} & Aligned subtraction & Unencoded \\
\texttt{05.05a} & Frobenius-norm invariance & Unencoded \\
\texttt{05.05b} & Operator-norm invariance & Unencoded \\
\texttt{05.06} & Equivariant path residual & Unencoded \\
\texttt{05.07} & Gauge--truncation counterexample & Unencoded \\
\texttt{05.08} & Shell-Frobenius invariance & Unencoded \\
\bottomrule
\end{tabular}
\end{center}

Every Isabelle and Rocq target is unencoded. Eight Lean targets are unencoded.
No contract is cross-checked. The package remains \texttt{proposed}, while
\texttt{PRF-05.01} is individually \texttt{in development} because the
multi-backend completion and review criteria have not been met.

\section{Novelty and evidence boundary}

The finite-dimensional projector, compression, pullback, subtraction, and norm
identities are standard linear algebra. The project does not claim those
identities as new mathematics. A possible contribution is their
application-specific organization into one prover-neutral contract set,
independent Lean--Isabelle--Rocq encodings, semantic conformance review, dual
common-space impurity extraction, and composition with later locality,
continuum, and compatibility arguments. Priority over prior formalization has
not been established.

Potentially substantive mathematical results lie in the assumptions and bounds
that connect the elementary layer to the physical reduction: existence and
regularity of an adequate identification, gauge-constrained locality,
quantitative excluded-space control, well-posed crossover criteria,
operator-to-observable bounds, certified model-class incompatibility, and
atomistic-to-continuum consistency. Every such result remains proposed.

Proof establishes only the stated conclusion from the stated assumptions.
Numerical verification must still test implementations and bound usefulness;
scientific validation must still test the physical model for a declared use;
and human review must still decide whether the assumptions and interpretation
are acceptable.
